\documentclass[10pt,a4paper,landscape]{article}
\usepackage[utf8]{inputenc}
\usepackage[T1]{fontenc}
\usepackage{lmodern}
\usepackage[margin=0.7cm]{geometry}
\usepackage{tikz}
\usetikzlibrary{calc,arrows.meta}
\usepackage{parskip}
\pagestyle{empty}
\renewcommand{\familydefault}{\sfdefault}

% ================= MAKRA SYMBOLI (IEC 60617, konwencja EPLAN) =================
\tikzset{every picture/.style={line width=0.55pt}, lbl/.style={font=\tiny, inner sep=0.5pt, fill=white},
tag/.style={font=\scriptsize\bfseries, inner sep=1pt}, zn/.style={font=\tiny, inner sep=0.5pt},
ods/.style={font=\tiny\itshape, inner sep=0.5pt}}

% zestyk zwierny pionowy; dolny zacisk (#1,#2), gorny (#1,#2+1.2)
\newcommand{\ZNO}[2]{\draw (#1,#2) -- (#1,{#2+0.35}) ({#1},{#2+0.35}) -- ({#1-0.3},{#2+1.02}) (#1,{#2+1.02}) -- (#1,{#2+1.2});}
% zestyk rozwierny pionowy
\newcommand{\ZNC}[2]{\draw (#1,#2) -- (#1,{#2+0.35}) ({#1},{#2+0.35}) -- ({#1-0.32},{#2+1.06}) (#1,{#2+0.78}) -- ({#1-0.22},{#2+0.78}) (#1,{#2+0.78}) -- (#1,{#2+1.2});}
% rozlacznik (zwierny + poprzeczka odlacznika na gornym zacisku)
\newcommand{\ROZL}[2]{\ZNO{#1}{#2}\draw ({#1-0.18},{#2+1.02}) -- ({#1+0.18},{#2+1.02});}
% bezpiecznik pionowy
\newcommand{\FUZ}[2]{\draw (#1,#2) -- (#1,{#2+1.2}) ({#1-0.18},{#2+0.22}) rectangle ({#1+0.18},{#2+0.98});}
% cewka pionowa: dolny zacisk (#1,#2), gorny (#1,#2+1.4)
\newcommand{\COILV}[2]{\draw (#1,#2) -- (#1,{#2+0.42}) ({#1-0.42},{#2+0.42}) rectangle ({#1+0.42},{#2+0.98}) (#1,{#2+0.98}) -- (#1,{#2+1.4});}
% przekladnik pradowy: okrag na pionowym przewodzie, centrum (#1,#2)
\newcommand{\CTs}[2]{\draw (#1,#2) circle (0.3); \draw ({#1+0.3},#2) -- ({#1+0.72},#2);}
% zacisk
\newcommand{\TERM}[2]{\draw[fill=white] (#1,#2) circle (0.07);}
% gniazdo (kielich otwarty w gore) + pin z gory
\newcommand{\SOCKup}[2]{\draw ({#1-0.22},{#2+0.1}) arc (180:360:0.22); \draw (#1,{#2+0.55}) -- (#1,{#2+0.06});}
% gniazdo (kielich otwarty w dol) + pin w dol
\newcommand{\SOCKdn}[2]{\draw ({#1-0.22},{#2-0.1}) arc (180:0:0.22); \draw (#1,{#2-0.55}) -- (#1,{#2-0.06});}
% naped E-STOP (grzybek) po lewej zestyku na wys. blade
\newcommand{\ESTOP}[2]{\draw[dashed] ({#1-0.16},{#2+0.7}) -- ({#1-0.6},{#2+0.7}); \draw ({#1-0.6},{#2+0.7}) -- ({#1-0.6},{#2+0.9}) ({#1-0.85},{#2+0.9}) -- ({#1-0.35},{#2+0.9}) ({#1-0.85},{#2+0.9}) arc (180:0:0.25);}
% naped krancowki (rolka)
\newcommand{\ROLL}[2]{\draw[dashed] ({#1-0.16},{#2+0.7}) -- ({#1-0.55},{#2+0.7}); \draw ({#1-0.55},{#2+0.55}) -- ({#1-0.55},{#2+0.85}); \draw ({#1-0.55},{#2+0.55}) circle (0.09);}
% naped termiczny theta
\newcommand{\THETA}[2]{\draw[dashed] ({#1-0.17},{#2+0.7}) -- ({#1-0.5},{#2+0.7}); \node[zn, left] at ({#1-0.5},{#2+0.7}) {$\theta$};}
% przycisk NO (START)
\newcommand{\PUSH}[2]{\draw[dashed] ({#1-0.16},{#2+0.7}) -- ({#1-0.55},{#2+0.7}); \draw ({#1-0.55},{#2+0.55}) -- ({#1-0.55},{#2+0.85}) ({#1-0.8},{#2+0.55}) -- ({#1-0.8},{#2+0.85}) ({#1-0.8},{#2+0.7}) -- ({#1-0.95},{#2+0.7});}
% silnik: kolo M, centrum, podpis
\newcommand{\MOT}[3]{\draw (#1,#2) circle (0.42); \node[font=\scriptsize] at (#1,{#2+0.08}) {M}; \node[font=\tiny] at (#1,{#2-0.16}) {#3};}
% strzalka potencjalu w prawo/lewo
\newcommand{\POTR}[2]{\draw (#1,#2) -- ({#1+0.35},{#2+0.12}) -- ({#1+0.35},{#2-0.12}) -- cycle;}

% ================= RAMA ARKUSZA EPLAN =================
% #1 oznaczenie strony, #2 opis strony, #3 nr strony, #4 liczba stron
\newcommand{\ramkaE}[4]{%
\draw[line width=0.8pt] (0,0) rectangle (28.0,18.8);
\draw[line width=0.5pt] (0.5,0.5) rectangle (27.5,18.3);
% kolumny 1..10 gora/dol
\foreach \i in {1,...,10}{
  \draw ({0.5+(\i-1)*2.7},18.3) -- ({0.5+(\i-1)*2.7},18.8) ({0.5+(\i-1)*2.7},0.0) -- ({0.5+(\i-1)*2.7},0.5);
  \node[font=\scriptsize] at ({0.5+(\i-1)*2.7+1.35},18.55) {\i};
  \node[font=\scriptsize] at ({0.5+(\i-1)*2.7+1.35},0.25) {\i};}
% wiersze A..F
\foreach \i/\l in {0/A,1/B,2/C,3/D,4/E,5/F}{
  \draw (0.0,{18.3-\i*2.966}) -- (0.5,{18.3-\i*2.966}) (27.5,{18.3-\i*2.966}) -- (28.0,{18.3-\i*2.966});
  \node[font=\scriptsize] at (0.25,{18.3-\i*2.966-1.48}) {\l};
  \node[font=\scriptsize] at (27.75,{18.3-\i*2.966-1.48}) {\l};}
% tabelka rysunkowa (prawy dol)
\draw[line width=0.5pt] (14.0,0.5) rectangle (27.5,2.6);
\draw (14.0,1.9) -- (27.5,1.9) (14.0,1.2) -- (23.2,1.2) (18.6,0.5) -- (18.6,2.6) (23.2,0.5) -- (23.2,2.6) (25.6,1.2) -- (25.6,1.9);
\node[font=\tiny, anchor=north west] at (14.05,2.55) {firma};
\node[font=\small\bfseries] at (16.3,2.2) {AMPERE POINT};
\node[font=\tiny, anchor=north west] at (14.05,1.85) {projekt};
\node[font=\scriptsize] at (16.3,1.5) {AMPERE\_POINT\_custom\_device};
\node[font=\tiny, anchor=north west] at (14.05,1.15) {nr projektu};
\node[font=\scriptsize] at (16.3,0.8) {AP-CD-001};
\node[font=\tiny, anchor=north west] at (18.65,2.55) {opis strony};
\node[font=\scriptsize, align=left, anchor=west] at (18.75,2.2) {#2};
\node[font=\tiny, anchor=north west] at (18.65,1.85) {data / edytor};
\node[font=\scriptsize, anchor=west] at (18.75,1.5) {03.07.2026 \; / \; D.~Cekała};
\node[font=\tiny, anchor=north west] at (18.65,1.15) {stan: wersja 2 robocza (do etapu 0)};
\node[font=\tiny, anchor=north west] at (23.25,2.55) {oznaczenie strony};
\node[font=\normalsize\bfseries] at (24.4,2.2) {#1};
\node[font=\tiny, anchor=north west] at (23.25,1.85) {norma};
\node[font=\scriptsize] at (24.4,1.5) {IEC 60617};
\node[font=\tiny, anchor=north west] at (25.65,1.85) {strona};
\node[font=\scriptsize] at (26.55,1.5) {#3 / #4};
\node[font=\scriptsize\bfseries, anchor=west] at (23.3,0.85) {custom device --- stacja testowa EVSE};
}

\begin{document}

% ============ STRONA TYTULOWA ============
\begin{center}
\vspace*{2.2cm}
{\LARGE\bfseries AMPERE POINT --- custom device}\\[6pt]
{\Large Schemat elektryczny --- wersja 2 (styl EPLAN)}\\[10pt]
{\normalsize 3 lipca 2026}\\[16pt]
\begin{minipage}{0.78\textwidth}\small
Arkusze E1--E3 w konwencji stron EPLAN: rama z siatką odniesień (kolumny 1--10, wiersze A--F),
tabelka rysunkowa, symbole IEC~60617, oznaczenia aparatów zgodne z wykazem wersji~1,
odsyłacze krzyżowe w formacie \emph{/arkusz.kolumna} (przy stykach --- adres cewki; pod cewkami --- lustro styków).
Treść obwodów identyczna z wersją~1 (topologia zatwierdzona Z1--Z6); wartości i typy aparatów robocze ---
do potwierdzenia po etapie~0. Dokument odwzorowuje docelowy wygląd projektu w EPLAN
(\texttt{AP-CD-001}), w którym powstanie wersja formalna.\\[8pt]
\textbf{E1} --- tor mocy 3F: X1 $\to$ Q0 $\to$ F1--F3 $\to$ B1 $\to$ T1--T3 $\to$ K1 $\to$ X2; odczep X3.\\
\textbf{E2} --- łańcuch bezpieczeństwa 24 V: S0, F11.1--12, S1, K2, K10, A2, S2/RB, cewki RB i K1.\\
\textbf{E3} --- moduł obciążenia (12 gałęzi K11--K22 / E1--E12), wentylatory M3/M4 z K2, sterowanie A1.
\end{minipage}
\end{center}
\newpage
% ============ ARKUSZ E1 — TOR MOCY ============
\begin{center}
\begin{tikzpicture}[x=1cm,y=1cm]
\ramkaE{=CA1+EAA/2}{E1 --- tor mocy 3F}{2}{4}
% ---- X1 gniazdo pojazdowe ----
\foreach \x/\n/\z in {5.0/L1/1, 7.2/L2/2, 9.4/L3/3, 11.6/N/4, 13.8/PE/5, 15.4/CP/6, 16.2/PP/7}{
  \draw (\x,17.65) -- (\x,17.25);
  \SOCKup{\x}{16.7}
  \node[zn] at (\x,17.85) {\n};
  \node[zn, right] at ({\x+0.1},16.55) {\z};
  \draw (\x,16.6) -- (\x,16.3);}
\draw[dashed] (4.3,16.25) rectangle (16.9,17.45);
\node[tag, left] at (4.25,17.3) {-X1};
\node[lbl, align=left] at (2.3,16.6) {gniazdo pojazdowe\\ Type 2 32 A 3F\\ (wtyk kabla DUT)};
\node[lbl, align=left] at (2.3,15.5) {blokada wtyku\\ -Y1 $\to$ /E3.3};
% ---- piony ----
\draw (5.0,16.3) -- (5.0,14.8) (7.2,16.3) -- (7.2,14.8) (9.4,16.3) -- (9.4,14.8) (11.6,16.3) -- (11.6,14.8);
\draw (13.8,16.3) -- (13.8,6.0); % PE
\node[lbl] at (14.15,9.0) {PE};
% ---- odczep X3 ----
\draw[fill] (5.0,15.9) circle (0.06); \draw (5.0,15.9) -- (17.0,15.9);
\draw[fill] (11.6,15.6) circle (0.06); \draw (11.6,15.6) -- (17.0,15.6);
\draw[fill] (13.8,15.3) circle (0.06); \draw (13.8,15.3) -- (17.0,15.3);
\draw (15.4,16.3) -- (15.4,15.0) -- (17.0,15.0);
\draw (16.2,16.3) -- (16.2,14.7) -- (17.0,14.7);
\foreach \y/\z in {15.9/1, 15.6/2, 15.3/3, 15.0/4, 14.7/5}{\TERM{17.0}{\y}\node[zn, right] at (17.1,\y) {\z};}
\draw[dashed] (16.75,14.45) rectangle (17.75,16.15);
\node[tag] at (17.25,16.4) {-X3};
\node[lbl, align=left, anchor=west] at (18.0,15.7) {odczep Type 2 $\to$ PM701E ($\le$1 m)\\ 1=L1, 2=N, 3=PE, 4=CP, 5=PP};
\node[ods, anchor=west] at (18.0,15.0) {SIGNAL OUTPUT (CP) $\to$ /E3.3};
\node[ods, anchor=west] at (18.0,14.6) {gniazda bananowe PM701E = przyłącze UT595};
% ---- Q0 rozlacznik 4P ----
\ROZL{5.0}{13.6}\ROZL{7.2}{13.6}\ROZL{9.4}{13.6}\ROZL{11.6}{13.6}
\draw[dashed] (4.55,14.25) -- (11.45,14.25);
\node[tag, left] at (4.2,14.2) {-Q0};
\node[lbl, left] at (4.2,13.8) {rozłącznik 40 A};
\foreach \x/\a/\b in {5.0/1/2, 7.2/3/4, 9.4/5/6, 11.6/7/8}{
  \node[zn, right] at ({\x+0.08},14.72) {\a}; \node[zn, right] at ({\x+0.08},13.7) {\b};}
% ---- F1-F3 ----
\draw (5.0,13.6) -- (5.0,13.1) (7.2,13.6) -- (7.2,13.1) (9.4,13.6) -- (9.4,13.1);
\draw (11.6,13.6) -- (11.6,11.4);
\FUZ{5.0}{11.9}\FUZ{7.2}{11.9}\FUZ{9.4}{11.9}
\node[tag, left] at (4.55,12.5) {-F1};\node[tag, left] at (6.75,12.5) {-F2};\node[tag, left] at (8.95,12.5) {-F3};
\node[lbl, left] at (4.2,12.05) {gG 32 A};
% ---- B1 okno CDSR ----
\draw (5.0,11.9) -- (5.0,9.9) (7.2,11.9) -- (7.2,9.9) (9.4,11.9) -- (9.4,9.9);
\draw (11.6,11.4) -- (11.6,9.9);
\draw[line width=1.0pt] (4.3,10.6) rectangle (12.3,11.35);
\node[tag, left] at (4.2,11.0) {-B1};
\node[lbl] at (8.3,10.32) {czujnik różnicowy AC/DC (CDSR 0.07-TP), okno L1+L2+L3+N};
\draw[-{Stealth}] (12.3,11.0) -- (13.1,11.0); \node[ods, right] at (13.1,11.0) {SPI /E3.3};
% ---- T1-T3 ----
\draw (5.0,9.9) -- (5.0,8.8) (7.2,9.9) -- (7.2,8.8) (9.4,9.9) -- (9.4,8.8) (11.6,9.9) -- (11.6,8.8);
\CTs{5.0}{9.35}\CTs{7.2}{9.35}\CTs{9.4}{9.35}
\node[tag, left] at (4.55,9.35) {-T1};\node[tag, left] at (6.75,9.35) {-T2};\node[tag, left] at (8.95,9.35) {-T3};
\node[ods] at (6.1,8.85) {/E3.3};\node[ods] at (8.3,8.85) {/E3.3};\node[ods] at (10.5,8.85) {/E3.3};
\node[lbl, left] at (4.2,9.0) {40 A / 1 V};
% ---- K1 stycznik glowny ----
\ZNO{5.0}{7.6}\ZNO{7.2}{7.6}\ZNO{9.4}{7.6}\ZNO{11.6}{7.6}
\draw[dashed] (4.55,8.25) -- (11.45,8.25);
\node[tag, left] at (4.2,8.35) {-K1};
\node[ods, left] at (4.2,7.95) {/E2.5};
\foreach \x/\a/\b in {5.0/1/2, 7.2/3/4, 9.4/5/6, 11.6/7/8}{
  \node[zn, right] at ({\x+0.08},8.72) {\a}; \node[zn, right] at ({\x+0.08},7.7) {\b};}
% ---- zejscie do X2 ----
\draw (5.0,7.6) -- (5.0,6.0) (7.2,7.6) -- (7.2,6.0) (9.4,7.6) -- (9.4,6.0) (11.6,7.6) -- (11.6,6.0);
\foreach \x/\n/\z in {5.0/L1/1, 7.2/L2/2, 9.4/L3/3, 11.6/N/4, 13.8/PE/5}{
  \SOCKdn{\x}{5.9}
  \draw (\x,5.35) -- (\x,5.0);
  \node[zn, right] at ({\x+0.1},6.05) {\z};
  \node[zn] at (\x,4.82) {\n};}
\draw[dashed] (4.3,5.15) rectangle (14.6,6.3);
\node[tag, left] at (4.25,6.15) {-X2};
\node[lbl, align=left] at (2.4,5.7) {gniazdo CEE 32 A 5p\\ $\to$ moduł B /E3.3};
\node[lbl] at (8.3,4.4) {przewód 5$\times$6 mm$^2$ H07RN-F};
\node[lbl, left] at (4.6,7.0) {6 mm$^2$};
% ---- uwagi ----
\node[lbl, align=left, anchor=north west] at (18.6,13.6) {\textbf{Uwagi:}\\ N bez zabezpieczenia --- rozłączany\\ czterobiegunowo (Q0, K1)\\ PE nieprzerwane od X1 do X2\\ B1: pomiar RCM, próg awaryjny $\to$ /E3.3\\ T1--T3: pomiar prądów $\to$ U2 /E3.3\\ K1: cewka i lustro styków $\to$ /E2.5};
\end{tikzpicture}
\end{center}
\newpage
% ============ ARKUSZ E2 — LANCUCH BEZPIECZENSTWA ============
\begin{center}
\begin{tikzpicture}[x=1cm,y=1cm]
\ramkaE{=CA1+EAA/3}{E2 --- łańcuch bezp. 24 V}{3}{4}
% ---- szyny ----
\draw (1.4,16.9) -- (14.8,16.9); \node[tag, left] at (1.35,16.9) {+24 V};
\POTR{14.8}{16.9} \node[ods, right] at (15.2,16.9) {z G1 /E3.5};
\draw (1.4,3.6) -- (14.8,3.6); \node[tag, left] at (1.35,3.6) {0 V};
\POTR{14.8}{3.6} \node[ods, right] at (15.2,3.6) {G1 /E3.5};
% ================= sciezka 1 (kol.1): S0 -> termiki -> S1 =================
\draw[fill] (3.0,16.9) circle (0.06);
\draw (3.0,16.9) -- (3.0,16.5);
\ZNC{3.0}{15.3}\ESTOP{3.0}{15.3}
\node[tag, right] at (3.2,15.9) {-S0};
\node[zn, right] at (3.08,16.35) {11}; \node[zn, right] at (3.08,15.4) {12};
\node[lbl, right] at (3.2,15.55) {E-STOP (panel A)};
\draw (3.0,15.3) -- (3.0,14.8);
\ZNC{3.0}{13.6}\THETA{3.0}{13.6}
\node[tag, right] at (3.2,14.2) {-F11.1};
\draw (3.0,13.6) -- (3.0,13.35);
\node[font=\scriptsize] at (3.0,13.05) {$\vdots$};
\draw (3.0,12.75) -- (3.0,12.7);
\ZNC{3.0}{11.5}\THETA{3.0}{11.5}
\node[tag, right] at (3.2,12.1) {-F11.12};
\draw[dashed] (1.95,11.25) rectangle (4.75,14.95);
\node[lbl, align=left, anchor=north west] at (1.95,11.1) {termiki KSD301 sekcji grzałek,\\ NC szeregowo; moduł B, złącze -X5 /E3.6};
\draw (3.0,11.5) -- (3.0,10.9);
\ZNC{3.0}{9.7}\ROLL{3.0}{9.7}
\node[tag, right] at (3.2,10.3) {-S1};
\node[lbl, right] at (3.2,9.95) {krańcówka pokrywy};
\node[zn, right] at (3.08,10.75) {11}; \node[zn, right] at (3.08,9.8) {12};
\draw (3.0,9.7) -- (3.0,8.9);
\draw[fill] (3.0,8.9) -- (2.85,8.55) -- (3.15,8.55) -- cycle;
\node[ods] at (3.0,8.25) {/E2.3};
% ================= sciezka 3 (kol.3): K2 -> K10 -> A2 -> [S2||RB] -> cewka RB =================
\draw[fill] (8.4,16.4) -- (8.25,16.05) -- (8.55,16.05) -- cycle;
\node[ods] at (8.4,16.6) {/E2.1};
\draw (8.4,16.05) -- (8.4,15.6);
\ZNO{8.4}{14.4}
\node[tag, left] at (8.1,15.0) {-K2};
\node[ods, left] at (8.1,14.6) {/E3.8};
\node[lbl, right] at (8.6,15.0) {potwierdzenie wentylatorów};
\node[zn, right] at (8.48,15.5) {13}; \node[zn, right] at (8.48,14.5) {14};
\draw (8.4,14.4) -- (8.4,13.9);
\ZNO{8.4}{12.7}
\node[tag, left] at (8.1,13.3) {-K10};
\node[ods, left] at (8.1,12.9) {/E3.5};
\node[lbl, right] at (8.6,13.3) {zezwolenie MCU};
\node[zn, right] at (8.48,13.8) {13}; \node[zn, right] at (8.48,12.8) {14};
\draw (8.4,12.7) -- (8.4,12.2);
\ZNO{8.4}{11.0}
\node[tag, left] at (8.1,11.6) {-A2};
\node[ods, left] at (8.1,11.2) {/E3.5};
\node[lbl, right] at (8.6,11.6) {watchdog impulsowy};
\draw (8.4,11.0) -- (8.4,10.5);
% galaz rownolegla
\draw[fill] (8.4,10.5) circle (0.06);
\draw (8.4,10.5) -- (10.6,10.5);
\PUSH{8.4}{8.9}\ZNO{8.4}{8.9}
\node[tag, anchor=east] at (7.3,10.15) {-S2};
\node[lbl, anchor=east] at (7.3,9.8) {START};
\node[zn, right] at (8.48,10.0) {13}; \node[zn, right] at (8.48,9.0) {14};
\draw (10.6,10.5) -- (10.6,10.1);
\ZNO{10.6}{8.9}
\node[tag, right] at (10.8,9.5) {-RB};
\node[ods, right] at (10.8,9.15) {13-14 \; /E2.3};
\draw (10.6,8.9) -- (10.6,8.5) -- (8.4,8.5);
\draw[fill] (8.4,8.5) circle (0.06);
\node[lbl, right] at (10.8,8.7) {samopodtrzymanie};
\draw (8.4,8.9) -- (8.4,8.0);
\COILV{8.4}{6.6}
\node[tag, left] at (7.9,7.3) {-RB};
\node[zn, right] at (8.48,7.85) {A1}; \node[zn, right] at (8.48,6.75) {A2};
\node[lbl, left] at (7.9,6.9) {przekaźnik bezpieczeństwa};
\draw (8.4,6.6) -- (8.4,3.6);
\draw[fill] (8.4,3.6) circle (0.06);
% lustro stykow RB
\node[ods, align=left, anchor=north west] at (8.7,6.1) {\textbf{-RB}\\ 13-14 \; /E2.3\\ 23-24 \; /E2.5};
% ================= sciezka 5 (kol.5): RB:23-24 -> cewka K1 =================
\draw[fill] (13.0,16.9) circle (0.06);
\draw (13.0,16.9) -- (13.0,15.8);
\ZNO{13.0}{14.6}
\node[tag, left] at (12.7,15.2) {-RB};
\node[ods, left] at (12.7,14.8) {/E2.3};
\node[zn, right] at (13.08,15.7) {23}; \node[zn, right] at (13.08,14.7) {24};
\node[lbl, right] at (13.2,15.2) {styk wymuszony mech.};
\draw (13.0,14.6) -- (13.0,13.2);
\COILV{13.0}{11.8}
\node[tag, left] at (12.5,12.5) {-K1};
\node[zn, right] at (13.08,13.05) {A1}; \node[zn, right] at (13.08,11.95) {A2};
\node[lbl, right] at (13.5,12.5) {stycznik główny /E1};
\draw (13.0,11.8) -- (13.0,3.6);
\draw[fill] (13.0,3.6) circle (0.06);
% lustro stykow K1
\node[ods, align=left, anchor=north west] at (13.3,10.9) {\textbf{-K1}\\ 1-2 \; /E1.2\\ 3-4 \; /E1.3\\ 5-6 \; /E1.4\\ 7-8 \; /E1.5};
% ---- adnotacje ----
\node[lbl, align=left, anchor=north west] at (16.5,14.0) {\textbf{Zasada:} tor 24 V cewki -K1 przechodzi wyłącznie\\ przez styki sprzętowe. Awaria członu, zanik 24 V,\\ zawieszenie firmware (zanik impulsów -A2)\\ lub E-STOP $\Rightarrow$ -RB odpada $\Rightarrow$ -K1 otwiera\\ tor mocy (/E1). Ponowne załączenie tylko -S2.};
\node[lbl, align=left, anchor=north west] at (16.5,11.2) {\textbf{Odczyty diagnostyczne (optoizolacja):}\\ stan -S0, -F11.x, -S1, -K2 $\to$ A1 /E3.4\\ (nie ingerują w tor łańcucha)};
\node[lbl, align=left, anchor=north west] at (16.5,9.3) {\textbf{Poza łańcuchem:} B1 (RCM) $\to$ zdjęcie -K10 /E3.5;\\ bezpieczniki gG /E1.2; rozłącznik -Q0 /E1.2;\\ blokada -Y1 wtyku Type 2 w stanie C /E3.3};
\end{tikzpicture}
\end{center}
\newpage
% ============ ARKUSZ E3 — MODUL OBCIAZENIA + STEROWANIE ============
\begin{center}
\begin{tikzpicture}[x=1cm,y=1cm]
\ramkaE{=CA1+EAA/4}{E3 --- obciążenie i sterowanie}{4}{4}
% ================= MODUL B: load bank =================
\draw[dashed] (7.5,10.9) rectangle (25.6,17.7);
\node[tag, anchor=south west] at (7.6,17.75) {MODUŁ B --- obciążenie (zasilanie z -X2 /E1.2--5, złącze sterownicze -X5)};
% szyny fazowe
\POTR{7.7}{16.9}\draw (8.05,16.9) -- (25.2,16.9); \node[ods, left] at (7.65,16.9) {L1 /E1.2};
\POTR{7.7}{16.5}\draw (8.05,16.5) -- (25.2,16.5); \node[ods, left] at (7.65,16.5) {L2 /E1.3};
\POTR{7.7}{16.1}\draw (8.05,16.1) -- (25.2,16.1); \node[ods, left] at (7.65,16.1) {L3 /E1.4};
\POTR{7.7}{11.5}\draw (8.05,11.5) -- (25.2,11.5); \node[ods, left] at (7.65,11.5) {N /E1.5};
% 12 galezi: K11-K14 (L1), K15-K18 (L2), K19-K22 (L3)
\foreach \x/\k/\e/\rail in {8.9/-K11/-E1/16.9, 10.3/-K12/-E2/16.9, 11.7/-K13/-E3/16.9, 13.1/-K14/-E4/16.9,
                            14.5/-K15/-E5/16.5, 15.9/-K16/-E6/16.5, 17.3/-K17/-E7/16.5, 18.7/-K18/-E8/16.5,
                            20.1/-K19/-E9/16.1, 21.5/-K20/-E10/16.1, 22.9/-K21/-E11/16.1, 24.3/-K22/-E12/16.1}{
  \draw[fill] (\x,\rail) circle (0.055);
  \draw (\x,\rail) -- (\x,15.1);
  \ZNO{\x}{13.9}
  \node[zn, left] at ({\x-0.3},14.75) {\k};
  \draw (\x,13.9) -- (\x,13.35);
  \draw ({\x-0.18},12.35) rectangle ({\x+0.18},13.35);
  \node[zn, left] at ({\x-0.22},12.85) {\e};
  \draw (\x,12.35) -- (\x,11.5);
  \draw[fill] (\x,11.5) circle (0.055);}
% moce stopni
\foreach \x/\p in {8.9/{1,0}, 10.3/{1,5}, 11.7/{2,0}, 13.1/{3,0}, 14.5/{1,0}, 15.9/{1,5}, 17.3/{2,0}, 18.7/{3,0}, 20.1/{1,0}, 21.5/{1,5}, 22.9/{2,0}, 24.3/{3,0}}{
  \node[zn] at (\x,12.1) {\p};}
\node[lbl] at (16.6,11.05) {grzałki żebrowane 230 V, moce w kW; termiki $\theta$ -F11.1\ldots12 na sekcjach $\to$ /E2.1; styki -K11\ldots-K22: cewki /E3.7--10};
% ================= wentylatory =================
\draw[dashed] (20.6,7.3) rectangle (26.4,10.3);
\node[tag, anchor=south west] at (20.7,10.35) {wentylatory modułu B};
\node[ods, left] at (21.0,9.6) {X4:L /E3.5};
\draw (21.05,9.6) -- (21.7,9.6);
\draw (21.7,9.25) rectangle (22.5,9.95); \node[font=\scriptsize] at (22.1,9.6) {I$>$};
\node[tag] at (22.1,10.15) {-K2};
\node[ods] at (22.1,8.95) {styk /E2.3};
\draw (22.5,9.6) -- (23.6,9.6);
\draw[fill] (23.6,9.6) circle (0.055);
\draw (23.6,9.6) -- (23.6,9.15); \MOT{23.6}{8.7}{1$\sim$}
\draw (23.6,9.6) -- (25.2,9.6) -- (25.2,9.15); \MOT{25.2}{8.7}{1$\sim$}
\node[zn, left] at (23.1,8.7) {-M3}; \node[zn, right] at (25.7,8.7) {-M4};
\draw (23.6,8.28) -- (23.6,7.9) -- (25.2,7.9) -- (25.2,8.28);
\draw (24.4,7.9) -- (24.4,7.6); \node[ods, below] at (24.4,7.6) {X4:N};
% ================= zasilacz G1 =================
\node[ods, left] at (1.3,10.0) {X4:1 (L)};
\node[ods, left] at (1.3,9.6) {X4:2 (N)};
\node[ods, left] at (1.3,9.2) {X4:3 (PE)};
\foreach \y in {10.0,9.6,9.2}{\TERM{1.45}{\y}\draw (1.52,\y) -- (2.0,\y);}
\draw (2.0,8.9) rectangle (4.6,10.3);
\node[font=\scriptsize\bfseries, align=center] at (3.3,9.85) {-G1};
\node[zn, align=center] at (3.3,9.35) {zasilacz\\ 24 / 5 / 3,3 V};
\draw[-{Stealth}] (4.6,9.95) -- (5.4,9.95); \node[ods, right] at (5.4,9.95) {24 V $\to$ /E2.1, cewki, -Y1};
\draw[-{Stealth}] (4.6,9.25) -- (5.4,9.25); \node[ods, right] at (5.4,9.25) {5 / 3,3 V $\to$ -A1, czujniki};
\node[lbl, align=left, anchor=north west] at (2.0,8.6) {zasilanie logiki z przewodu Schuko\\ (nie z badanej ładowarki)};
% ================= A1 ESP32 =================
\draw (1.5,3.4) rectangle (7.0,7.6);
\node[font=\scriptsize\bfseries] at (4.25,7.25) {-A1 \; ESP32-S3 --- płyta logiki};
\foreach \y/\p/\t in {6.7/SPI/{-B1 /E1.3; -U2 ATM90E36A $\leftarrow$ -T1\ldots-T3 /E1.3},
                      6.15/I2C/{-B2 MLX90640; -B3, -B4 MLX90614; -U3, -U4 MCP23017},
                      5.6/UART/{-P1 HMI DWIN 7'' (START, LED, buzzer)},
                      5.05/{1-W}/{-B5.1\ldots8 DS18B20 (szyny, K1, komora, otoczenie)},
                      4.5/ADC/{CP, PP z -X3 /E1.6 (dzielnik + ograniczniki)},
                      3.95/DI/{odczyty łańcucha (opto) /E2.1\\ -S0, -F11.x, -S1, -K2}}{
  \TERM{7.0}{\y}
  \node[zn, left] at (6.9,\y) {\p};
  \draw (7.07,\y) -- (7.6,\y);
  \node[ods, right, align=left] at (7.6,\y) {\t};}
% DO: cewki K10 i A2 (przy A1) + magistrala do ULN
\TERM{4.25}{3.4}
\node[zn, above right] at (4.3,3.42) {DO};
\draw (4.25,3.33) -- (4.25,2.2) -- (13.65,2.2);
\node[lbl] at (5.6,1.95) {DO (opto, low-side)};
\draw[fill] (8.2,2.2) circle (0.05);
\draw[fill] (12.4,2.2) circle (0.05);
\draw (12.4,2.2) -- (12.4,3.3); \COILV{12.4}{3.3}
\node[tag, right] at (12.9,4.35) {-K10};
\node[ods, right] at (12.9,4.0) {styk /E2.3};
\draw (12.4,4.7) -- (12.4,4.95); \node[ods, above] at (12.4,4.95) {24 V};
\draw (13.65,2.2) -- (13.65,3.3);
\draw (13.25,3.3) rectangle (14.05,4.1); \node[zn, align=center] at (13.65,3.7) {-A2\\ WDT};
\node[ods, anchor=east] at (13.2,3.0) {styk /E2.3};
\draw (13.65,4.1) -- (13.65,4.35); \node[ods, anchor=east] at (13.55,4.55) {24 V};
\node[lbl] at (10.4,2.85) {impulsy watchdog + zezwolenie};
% ================= A3 generator uplywu =================
\draw (8.4,8.6) rectangle (13.6,10.3);
\node[font=\scriptsize\bfseries] at (11.0,9.95) {-A3 --- generator upływu};
\node[zn, align=left, anchor=north west] at (8.55,9.7) {DC 0--20 mA / AC (rampa z DAC);\\ pomiar: bocznik + AMC1301 $\to$ -A1;\\ przyłącza: L1 przez -K30 /E1.2, PE};
\node[ods, anchor=north west] at (8.55,8.5) {test RDC-DD 6 mA DC i RCD ładowarki};
% ================= Y1 =================
\draw (8.2,2.2) -- (8.2,2.3); \COILV{8.2}{2.3}
\node[tag, right] at (8.7,3.3) {-Y1};
\node[ods, right] at (8.7,2.95) {blokada wtyku -X1 /E1.6};
\draw (8.2,3.7) -- (8.2,3.95); \node[ods, above] at (8.2,3.95) {24 V};
% ================= rzad cewek stopni =================
\draw (14.4,6.4) -- (26.3,6.4); \node[ods, left] at (14.35,6.4) {24 V};
\foreach \i/\x in {11/14.9, 12/15.6, 13/16.3, 14/17.0, 15/17.7, 16/18.4, 17/19.1, 18/19.8, 19/20.5, 20/21.2, 21/21.9, 22/22.6}{
  \draw[fill] (\x,6.4) circle (0.05);
  \draw (\x,6.4) -- (\x,6.15);
  \draw ({\x-0.28},5.55) rectangle ({\x+0.28},6.15);
  \draw (\x,5.55) -- (\x,5.3);
  \node[zn, rotate=90, anchor=east] at (\x,5.2) {-K\i};}
\foreach \i/\x in {30/23.3, 31/24.0, 32/24.7, 33/25.4}{
  \draw[fill] (\x,6.4) circle (0.05);
  \draw (\x,6.4) -- (\x,6.15);
  \draw ({\x-0.28},5.55) rectangle ({\x+0.28},6.15);
  \draw (\x,5.55) -- (\x,5.3);
  \node[zn, rotate=90, anchor=east] at (\x,5.2) {-K\i};}
\draw (14.4,4.35) -- (26.3,4.35); \node[ods, left] at (14.35,4.35) {0 V};
\foreach \x in {14.9,15.6,16.3,17.0,17.7,18.4,19.1,19.8,20.5,21.2,21.9,22.6,23.3,24.0,24.7,25.4}{
  \draw (\x,5.3) -- (\x,4.35); \draw[fill] (\x,4.35) circle (0.05);}
\draw (13.0,6.9) rectangle (17.4,7.7); \node[zn, align=center] at (15.2,7.3) {-U5/-U6 ULN2803 $\leftarrow$ -U3/-U4 (I2C -A1)};
\draw[-{Stealth}] (15.2,6.9) -- (15.2,6.55);
\node[lbl, align=left, anchor=north west] at (14.4,3.95) {sterowanie stopniami: styki -K11\ldots-K22 /E3.3--10 (góra arkusza);\\ -K30\ldots-K33: separacja -A3, -U2, dzielników od L/N/PE (tryb izolacji /E1.3)};
\end{tikzpicture}
\end{center}

\end{document}
